% problems3.tex — Problem Set and Answers for Chapter 3

\begin{problemset}

\noindent\textsc{Analytical Exercises}
\bigskip

\exercise[1]
Prove: A symmetric and idempotent matrix is positive semidefinite.

\exercise[2]
We wish to prove Proposition~3.4.  To avoid inessential complications, consider
the case where $K = 1$ and $L = 1$ so that $x_i$ and $z_i$ are scalars.  So (3.5.5)
becomes
\[
  \hat{\varepsilon}_i^2 = \varepsilon_i^2
  - 2(\hat{\delta} - \delta)\, z_i\, \varepsilon_i
  + (\hat{\delta} - \delta)^2 z_i^2,
\]
and the $\hat{\mathbf{S}}$ in (3.5.10) becomes
\[
  \frac{1}{n} \sum_{i=1}^{n} x_i^2\, \varepsilon_i^2
  - 2(\hat{\delta} - \delta) \frac{1}{n} \sum_{i=1}^{n} z_i\, x_i^2\, \varepsilon_i
  + (\hat{\delta} - \delta)^2 \frac{1}{n} \sum_{i=1}^{n} z_i^2\, x_i^2.
  \tag{$*$}
\]

\begin{enumerate}[label=(\alph*)]
\item Show: The first term on the RHS of $(*)$ converges in probability to $S$
  $(\equiv \E(x_i^2\, \varepsilon_i^2))$.

\item Use the Cauchy--Schwarz inequality
  \[
    \E(|X \cdot Y|) \leq \sqrt{\E(X^2) \cdot \E(Y^2)}
  \]
  to show that $\E(z_i x_i^2 \varepsilon_i)$ exists and is finite.
  \textbf{Hint:} $z_i x_i^2 \varepsilon_i$ is the product of $x_i \varepsilon_i$ and $x_i z_i$.
  Fact: $\E(x)$ exists and is finite if and only if $\E(|x|) < \infty$.
  Where do we use Assumption~3.6?

\item Show that the second term on the RHS of $(*)$ converges to zero in probability.

\item In a similar fashion, show that the third term on the RHS of $(*)$ converges
  to zero in probability.
\end{enumerate}

\exercise[3]
We wish to prove the algebraic result (3.5.11).  What needs to be proved is that
$\mathbf{A} - \mathbf{B}$ is positive semidefinite where
\[
  \mathbf{A} \equiv
  (\boldsymbol{\Sigma}_{\mathbf{xz}}' \mathbf{W} \boldsymbol{\Sigma}_{\mathbf{xz}})^{-1}\,
  \boldsymbol{\Sigma}_{\mathbf{xz}}' \mathbf{W} \mathbf{S} \mathbf{W}
  \boldsymbol{\Sigma}_{\mathbf{xz}}\,
  (\boldsymbol{\Sigma}_{\mathbf{xz}}' \mathbf{W} \boldsymbol{\Sigma}_{\mathbf{xz}})^{-1}
\]
and
\[
  \mathbf{B} \equiv
  (\boldsymbol{\Sigma}_{\mathbf{xz}}' \mathbf{S}^{-1}
  \boldsymbol{\Sigma}_{\mathbf{xz}})^{-1}.
\]

A result from matrix algebra states:

Let $\mathbf{A}$ and $\mathbf{B}$ be two positive definite and symmetric matrices.
$\mathbf{A} - \mathbf{B}$ is positive semidefinite if and only if
$\mathbf{B}^{-1} - \mathbf{A}^{-1}$ is positive semidefinite.

Therefore, what needs to be proved is that
\[
  \mathbf{Q} \equiv
  \boldsymbol{\Sigma}_{\mathbf{xz}}' \mathbf{S}^{-1} \boldsymbol{\Sigma}_{\mathbf{xz}}
  - \boldsymbol{\Sigma}_{\mathbf{xz}}' \mathbf{W} \boldsymbol{\Sigma}_{\mathbf{xz}}\,
  (\boldsymbol{\Sigma}_{\mathbf{xz}}' \mathbf{W} \mathbf{S} \mathbf{W}
  \boldsymbol{\Sigma}_{\mathbf{xz}})^{-1}\,
  \boldsymbol{\Sigma}_{\mathbf{xz}}' \mathbf{W} \boldsymbol{\Sigma}_{\mathbf{xz}}
\]
is positive semidefinite.

\begin{enumerate}[label=(\alph*)]
\item Since $\mathbf{S}$ ($K \times K$) is positive definite, there exists a nonsingular $K \times K$
  matrix $\mathbf{C}$ such that $\mathbf{C}'\mathbf{C} = \mathbf{S}^{-1}$.  Define
  $\mathbf{A} \equiv \mathbf{C}\boldsymbol{\Sigma}_{\mathbf{xz}}$.  Show that
  $\hat{\mathbf{B}} \hat{\mathbf{S}}^{-1} \hat{\mathbf{B}} = \mathbf{C}'\mathbf{M}_G\mathbf{C}$,
  where
  \[
    \mathbf{H} = \mathbf{C}\boldsymbol{\Sigma}_{\mathbf{xz}}, \quad
    \mathbf{M}_G = \mathbf{I}_K - \mathbf{G}(\mathbf{G}'\mathbf{G})^{-1}\mathbf{G}',
    \quad \mathbf{G} = \mathbf{C}'^{-1}\mathbf{W}\boldsymbol{\Sigma}_{\mathbf{xz}}.
  \]
  \textbf{Hint:} $\mathbf{C}^{-1}\mathbf{C}'^{-1} = \mathbf{S}$.

\item Show that $\mathbf{Q}$ is positive semidefinite.
  \textbf{Hint:} First show that $\mathbf{M}_G$ is symmetric and idempotent.  As you showed in
  Exercise~1, a symmetric and idempotent matrix is positive semidefinite.
\end{enumerate}

\exercise[4]
Suppose $\mathbf{z}_i$ is a strict subset of $\mathbf{x}_i$.  Which one is smaller in the matrix sense,
the asymptotic variance of the OLS estimator or the asymptotic variance of
the efficient two-step GMM estimator?  \textbf{Hint:} Without loss of generality, we can
assume that the first $L$ elements of the $K$-dimensional vector $\mathbf{x}_i$ are $\mathbf{z}_i$.
If the $K \times K$ matrix $\mathbf{W}$ is given by
\[
  \mathbf{W} = \begin{bmatrix} \mathbf{A}^{-1} & \mathbf{0} \\
  \mathbf{0} & \mathbf{0} \end{bmatrix},
\]
where
\[
  \underset{(L \times L)}{\mathbf{A}} \equiv \E(\varepsilon_i^2\, \mathbf{z}_i \mathbf{z}_i')
  = \text{leading } L \times L \text{ submatrix of } \mathbf{S},
\]
then the asymptotic variance of the OLS estimator can be written as (3.5.1).
If you did Exercise~3, you can verify that the algebraic relation (3.5.11) does not
require $\mathbf{W}$ to be positive definite, as long as
$\boldsymbol{\Sigma}_{\mathbf{xz}}' \mathbf{W} \mathbf{S} \mathbf{W} \boldsymbol{\Sigma}_{\mathbf{xz}}$
is nonsingular.

\exercise[5]
(optional) Prove Proposition~3.6 by following the steps below.

\begin{enumerate}[label=(\alph*)]
\item Show: $\mathbf{g}_n(\hat{\boldsymbol{\delta}}(\hat{\mathbf{S}}^{-1}))
  = \hat{\mathbf{B}}\, \bar{\mathbf{g}}$ where
  $\hat{\mathbf{B}} \equiv \mathbf{I}_K
  - \mathbf{S}_{\mathbf{xz}}(\mathbf{S}_{\mathbf{xz}}'
  \hat{\mathbf{S}}^{-1} \mathbf{S}_{\mathbf{xz}})^{-1}
  \mathbf{S}_{\mathbf{xz}}' \hat{\mathbf{S}}^{-1}$.

\item Since $\hat{\mathbf{S}}$ is positive definite, there exists a nonsingular $K \times K$ matrix
  $\mathbf{C}$ such that $\mathbf{C}'\mathbf{C} = \hat{\mathbf{S}}^{-1}$.  Define
  $\mathbf{A} \equiv \mathbf{C}\mathbf{S}_{\mathbf{xz}}$.  Show that
  $\hat{\mathbf{B}}\hat{\mathbf{S}}^{-1}\hat{\mathbf{B}} = \mathbf{C}'\mathbf{M}\mathbf{C}$,
  where $\mathbf{M} \equiv \mathbf{I}_K - \mathbf{A}(\mathbf{A}'\mathbf{A})^{-1}\mathbf{A}'$.
  What is the rank of $\mathbf{M}$?  \textbf{Hint:} The rank of an idempotent matrix equals its trace.

\item Let $\mathbf{v} \equiv \sqrt{n}\, \mathbf{C}\, \bar{\mathbf{g}}$.
  Show that $\mathbf{v} \dto N(\mathbf{0}, \mathbf{I}_K)$.

\item Show that $J(\hat{\boldsymbol{\delta}}(\hat{\mathbf{S}}^{-1}), \hat{\mathbf{S}}^{-1})
  = \mathbf{v}'\mathbf{M}\mathbf{v}$.  What is its asymptotic distribution?
\end{enumerate}

\exercise[6]
Show that the $J$ statistic in Proposition~3.6 is numerically equal to
\[
  n \cdot \mathbf{s}_{\mathbf{xy}}' \hat{\mathbf{B}} \hat{\mathbf{S}}^{-1}
  \hat{\mathbf{B}} \mathbf{s}_{\mathbf{xy}},
\]
where $\hat{\mathbf{B}}$ is defined in Exercise~5(a).  \textbf{Hint:} From Exercise~5,
$J = n \cdot \bar{\mathbf{g}}' \hat{\mathbf{B}} \hat{\mathbf{S}}^{-1} \hat{\mathbf{B}} \bar{\mathbf{g}}$.
Show that $\hat{\mathbf{B}} \bar{\mathbf{g}} = \hat{\mathbf{B}} \mathbf{s}_{\mathbf{xy}}$.

\exercise[7]
(optional) Prove Proposition~3.7 by taking the following steps.  The notation
not introduced in this exercise is the same as in Exercise~5.

\begin{enumerate}[label=(\alph*)]
\item Prove:
  \[
    \mathbf{g}_{1n}(\tilde{\boldsymbol{\delta}})
    = \hat{\mathbf{B}}_1 \bar{\mathbf{g}}_1
    \quad \text{and} \quad
    \hat{\mathbf{B}}_1'(\hat{\mathbf{S}}_{11})^{-1} \hat{\mathbf{B}}_1
    = \mathbf{C}_1' \mathbf{M}_1 \mathbf{C}_1,
  \]
  where
  \begin{align*}
    \bar{\mathbf{g}}_1 &\equiv \frac{1}{n} \sum \mathbf{x}_{i1} \cdot \varepsilon_i, \\
    \hat{\mathbf{B}}_1 &\equiv \mathbf{I}_{K_1}
    - \mathbf{S}_{\mathbf{x}_1 \mathbf{z}}
    [\mathbf{S}_{\mathbf{x}_1 \mathbf{z}}'(\hat{\mathbf{S}}_{11})^{-1}
    \mathbf{S}_{\mathbf{x}_1 \mathbf{z}}]^{-1}
    \mathbf{S}_{\mathbf{x}_1 \mathbf{z}}'(\hat{\mathbf{S}}_{11})^{-1}, \\
    \mathbf{M}_1 &\equiv \mathbf{I}_{K_1}
    - \mathbf{A}_1(\mathbf{A}_1'\mathbf{A}_1)^{-1}\mathbf{A}_1', \\
    \mathbf{A}_1 &\equiv \mathbf{C}_1 \mathbf{S}_{\mathbf{x}_1 \mathbf{z}}.
  \end{align*}

\item Prove: $\rank(\mathbf{M}) = K - L$, $\rank(\mathbf{M}_1) = K_1 - L$.

\item Prove: You have proved in the previous exercise that
  $J = \mathbf{v}'\mathbf{M}\mathbf{v}$, where $\mathbf{v} \equiv \sqrt{n}\, \mathbf{C}\, \bar{\mathbf{g}}$.
  Here, prove: $J_1 = \mathbf{v}_1' \mathbf{M}_1 \mathbf{v}_1$, where
  $\mathbf{v}_1 \equiv \sqrt{n}\, \mathbf{C}_1 \bar{\mathbf{g}}_1$.

  To proceed further, define the $K \times K_1$ matrix $\mathbf{F}$ as
  \[
    \underset{(K \times K_1)}{\mathbf{F}}
    = \begin{bmatrix} \mathbf{I}_{K_1} \\ \mathbf{0} \end{bmatrix}
    \quad K_1 \text{ rows,} \quad K - K_1 \text{ rows.}
  \]
  Then: $\mathbf{x}_{i1} = \mathbf{F}'\mathbf{x}_i$,
  $\mathbf{S}_{\mathbf{x}_1 \mathbf{z}} = \mathbf{F}'\mathbf{S}_{\mathbf{xz}}$,
  $\bar{\mathbf{g}}_1 = \mathbf{F}'\bar{\mathbf{g}}$.

\item Prove $J - J_1 = \mathbf{v}'(\mathbf{M} - \mathbf{D})\mathbf{v}$ where
  $\mathbf{D} = \mathbf{C}'^{-1}\mathbf{F}\mathbf{C}_1'\mathbf{M}_1\mathbf{C}_1
  \mathbf{F}'\mathbf{C}^{-1}$.

\item Prove that $\mathbf{D}$ is symmetric and idempotent, with rank $K_1 - L$.
  \textbf{Hint:}
  $\mathbf{F}'\mathbf{C}^{-1}\mathbf{C}'^{-1}\mathbf{F}
  = \mathbf{F}'\hat{\mathbf{S}}\mathbf{F} = \hat{\mathbf{S}}_{11}$,
  $\mathbf{C}_1 \hat{\mathbf{S}}_{11} \mathbf{C}_1'
  = \mathbf{C}_1(\mathbf{C}_1'\mathbf{C}_1)^{-1}\mathbf{C}_1' = \mathbf{I}_{K_1}$.

\item Prove: $\mathbf{A}'\mathbf{D} = \mathbf{0}$.  \textbf{Hint:}
  \[
    \mathbf{A}'\mathbf{D}
    = (\mathbf{S}_{\mathbf{xz}}'\mathbf{C}')
    (\mathbf{C}'^{-1}\mathbf{F}\mathbf{C}_1'\mathbf{M}_1
    \mathbf{C}_1\mathbf{F}'\mathbf{C}^{-1})
    = (\mathbf{S}_{\mathbf{xz}}'\mathbf{F}\mathbf{C}_1'\mathbf{M}_1)
    (\mathbf{C}_1\mathbf{F}'\mathbf{C}^{-1}),
  \]
  so $\mathbf{S}_{\mathbf{xz}}'\mathbf{F}\mathbf{C}_1'\mathbf{M}_1
  = \mathbf{A}_1'\mathbf{M}_1$.

\item Prove: $\mathbf{M} - \mathbf{D}$ is symmetric and idempotent, with rank $K - K_1$.

\item Prove the desired result: $J - J_1 \dto \chi^2(K - K_1)$.

\item Step~(d) has established that $C = n \cdot \bar{\mathbf{g}}'
  \mathbf{C}'(\mathbf{M} - \mathbf{D})\mathbf{C}\, \bar{\mathbf{g}}$.  Show that $C$ can be
  written as
  \[
    n \cdot \mathbf{s}_{\mathbf{xy}}' \mathbf{C}'(\mathbf{M} - \mathbf{D})
    \mathbf{C}\, \mathbf{s}_{\mathbf{xy}}.
  \]
  \textbf{Hint:} $\mathbf{g}_n(\hat{\boldsymbol{\delta}})
  = \hat{\mathbf{B}}\mathbf{s}_{\mathbf{xy}}$.  So the numerical value of $C$ would have been
  the same if we replaced $\bar{\mathbf{g}}$ by $\mathbf{s}_{\mathbf{xy}}$ from the beginning in~(a).

\item Show that $\mathbf{M} - \mathbf{D}$ is positive semidefinite.  Thus, the quadratic form $C$ is
  nonnegative for any sample.
\end{enumerate}

\exercise[8]
(Optional, proof of numerical equivalence of Proposition~3.8)  In the restricted
efficient GMM, the $J$ function is minimized subject to the restriction
$\mathbf{R}\tilde{\boldsymbol{\delta}} = \mathbf{r}$.  For the time being, do not restrict $\hat{\mathbf{W}}$
to be equal to $\hat{\mathbf{S}}^{-1}$ and form the Lagrangian
\[
  \mathcal{L} = n \cdot (\mathbf{s}_{\mathbf{xy}}
  - \mathbf{S}_{\mathbf{xz}} \tilde{\boldsymbol{\delta}})'\,
  \hat{\mathbf{W}}\,
  (\mathbf{s}_{\mathbf{xy}} - \mathbf{S}_{\mathbf{xz}} \tilde{\boldsymbol{\delta}})
  + \boldsymbol{\lambda}'(\mathbf{R}\tilde{\boldsymbol{\delta}} - \mathbf{r}),
\]
where $\boldsymbol{\lambda}$ is the $\#\mathbf{r}$-dimensional Lagrange multipliers (recall: $\mathbf{r}$
is $\#\mathbf{r} \times 1$, $\mathbf{R}$ is $\#\mathbf{r} \times L$, and $\tilde{\boldsymbol{\delta}}$
is $L \times 1$).  Let $\tilde{\boldsymbol{\delta}}$ be the solution to the constrained minimization
problem.  (If $\hat{\mathbf{W}} = \hat{\mathbf{S}}^{-1}$, it is the restricted efficient GMM estimator
of $\boldsymbol{\delta}$.)

\begin{enumerate}[label=(\alph*)]
\item Let $\hat{\boldsymbol{\delta}}(\hat{\mathbf{W}})$ be the unrestricted GMM estimator that minimizes
  $J(\tilde{\boldsymbol{\delta}}, \hat{\mathbf{W}})$.  Show:
  \begin{align*}
    \tilde{\boldsymbol{\delta}}
    &= \hat{\boldsymbol{\delta}}(\hat{\mathbf{W}})
    - (\mathbf{S}_{\mathbf{xz}}' \hat{\mathbf{W}} \mathbf{S}_{\mathbf{xz}})^{-1}
    \mathbf{R}'\,
    [\mathbf{R}(\mathbf{S}_{\mathbf{xz}}' \hat{\mathbf{W}} \mathbf{S}_{\mathbf{xz}})^{-1}
    \mathbf{R}']^{-1}\,
    [\mathbf{R}\hat{\boldsymbol{\delta}}(\hat{\mathbf{W}}) - \mathbf{r}], \\
    \boldsymbol{\lambda}
    &= 2n \cdot
    [\mathbf{R}(\mathbf{S}_{\mathbf{xz}}' \hat{\mathbf{W}} \mathbf{S}_{\mathbf{xz}})^{-1}
    \mathbf{R}']^{-1}\,
    [\mathbf{R}\hat{\boldsymbol{\delta}}(\hat{\mathbf{W}}) - \mathbf{r}].
  \end{align*}
  \textbf{Hint:} The first-order conditions are
  \[
    -2n\, \mathbf{S}_{\mathbf{xz}}' \hat{\mathbf{W}} \mathbf{s}_{\mathbf{xy}}
    + 2n\, (\mathbf{S}_{\mathbf{xz}}' \hat{\mathbf{W}} \mathbf{S}_{\mathbf{xz}})
    \tilde{\boldsymbol{\delta}} + \mathbf{R}'\boldsymbol{\lambda} = \mathbf{0}.
  \]
  Combine this with the constraint $\mathbf{R}\tilde{\boldsymbol{\delta}} = \mathbf{r}$ to solve for
  $\boldsymbol{\lambda}$ and $\tilde{\boldsymbol{\delta}}$.

\item Show:
  \[
    J(\tilde{\boldsymbol{\delta}}, \hat{\mathbf{W}})
    = J(\hat{\boldsymbol{\delta}}(\hat{\mathbf{W}}), \hat{\mathbf{W}})
    + n \cdot (\hat{\boldsymbol{\delta}}(\hat{\mathbf{W}}) - \tilde{\boldsymbol{\delta}})'\,
    (\mathbf{S}_{\mathbf{xz}}' \hat{\mathbf{W}} \mathbf{S}_{\mathbf{xz}})\,
    (\hat{\boldsymbol{\delta}}(\hat{\mathbf{W}}) - \tilde{\boldsymbol{\delta}}).
  \]

\item Now set $\hat{\mathbf{W}} = \hat{\mathbf{S}}^{-1}$ and show that the Wald statistic $W$
  defined in (3.5.8) is numerically equal to $LR$ in (3.7.2).
\end{enumerate}

\exercise[9]
(GMM Hausman principle) Let
$\hat{\boldsymbol{\delta}}_1 \equiv \hat{\boldsymbol{\delta}}(\hat{\mathbf{W}}_1)$ and
$\hat{\boldsymbol{\delta}}_2 \equiv \hat{\boldsymbol{\delta}}(\hat{\mathbf{W}}_2)$ be two GMM
estimators with two different choices of the weighting matrix, $\hat{\mathbf{W}}_1$ and
$\hat{\mathbf{W}}_2$.

\begin{enumerate}[label=(\alph*)]
\item Show:
  \[
    \begin{bmatrix} \sqrt{n}(\hat{\boldsymbol{\delta}}_1 - \boldsymbol{\delta}) \\
    \sqrt{n}(\hat{\boldsymbol{\delta}}_2 - \boldsymbol{\delta}) \end{bmatrix}
    \dto N\!\left(\mathbf{0},\,
    \begin{bmatrix} \mathbf{A}_{11} & \mathbf{A}_{12} \\
    \mathbf{A}_{21} & \mathbf{A}_{22} \end{bmatrix}\right),
  \]
  where
  \begin{gather*}
    \mathbf{A}_{11} = \mathbf{Q}_1^{-1} \boldsymbol{\Sigma}_{\mathbf{xz}}'
    \mathbf{W}_1 \mathbf{S} \mathbf{W}_1 \boldsymbol{\Sigma}_{\mathbf{xz}}
    \mathbf{Q}_1^{-1}, \quad
    \mathbf{A}_{12} = \mathbf{Q}_1^{-1} \boldsymbol{\Sigma}_{\mathbf{xz}}'
    \mathbf{W}_1 \mathbf{S} \mathbf{W}_2 \boldsymbol{\Sigma}_{\mathbf{xz}}
    \mathbf{Q}_2^{-1}, \\
    \mathbf{A}_{21} = \mathbf{Q}_2^{-1} \boldsymbol{\Sigma}_{\mathbf{xz}}'
    \mathbf{W}_2 \mathbf{S} \mathbf{W}_1 \boldsymbol{\Sigma}_{\mathbf{xz}}
    \mathbf{Q}_1^{-1}, \quad
    \mathbf{A}_{22} = \mathbf{Q}_2^{-1} \boldsymbol{\Sigma}_{\mathbf{xz}}'
    \mathbf{W}_2 \mathbf{S} \mathbf{W}_2 \boldsymbol{\Sigma}_{\mathbf{xz}}
    \mathbf{Q}_2^{-1}, \\
    \mathbf{Q}_1 = \boldsymbol{\Sigma}_{\mathbf{xz}}' \mathbf{W}_1
    \boldsymbol{\Sigma}_{\mathbf{xz}}, \quad
    \mathbf{Q}_2 = \boldsymbol{\Sigma}_{\mathbf{xz}}' \mathbf{W}_2
    \boldsymbol{\Sigma}_{\mathbf{xz}}, \quad
    \plim_{n\to\infty} \hat{\mathbf{W}}_j = \mathbf{W}_j \; (j = 1, 2).
  \end{gather*}

\item Let $\mathbf{q} \equiv \hat{\boldsymbol{\delta}}_1 - \hat{\boldsymbol{\delta}}_2$.
  Show that $\sqrt{n}\, \mathbf{q} \dto N(\mathbf{0}, \avar(\mathbf{q}))$, where
  \[
    \avar(\mathbf{q})
    = \mathbf{A}_{11} + \mathbf{A}_{22} - \mathbf{A}_{12} - \mathbf{A}_{21}.
  \]

\item Set $\hat{\mathbf{W}}_2 = \hat{\mathbf{S}}^{-1}$ so that $\hat{\boldsymbol{\delta}}_2$ is efficient GMM.
  Show that
  \[
    \avar(\mathbf{q})
    = \mathbf{A}_{11}
    - (\boldsymbol{\Sigma}_{\mathbf{xz}}' \mathbf{S}^{-1}
    \boldsymbol{\Sigma}_{\mathbf{xz}})^{-1}
    = \avar(\hat{\boldsymbol{\delta}}(\hat{\mathbf{W}}_1))
    - \avar(\hat{\boldsymbol{\delta}}(\hat{\mathbf{S}}^{-1})).
  \]
\end{enumerate}

\exercise[10]
(Weak instruments.  This question is due to M.\ Watson.)  Consider the following
model:
\[
  y_i = z_i \delta + \varepsilon_i \quad (i = 1, 2, \ldots, n),
\]
where $y_i$, $z_i$, and $\varepsilon_i$ are scalar random variables.  Let $x_i$ denote a scalar
instrumental variable that is related to $z_i$ via the equation:
\[
  z_i = x_i \beta + v_i.
\]
Let
\[
  \boldsymbol{\eta}_i \equiv \begin{bmatrix} \varepsilon_i \\ v_i \end{bmatrix}, \quad
  \mathbf{g}_i \equiv \boldsymbol{\eta}_i \cdot x_i
  = \begin{bmatrix} g_{1i} \\ g_{2i} \end{bmatrix}
  = \begin{bmatrix} x_i \varepsilon_i \\ x_i v_i \end{bmatrix}.
\]
Assume that (1)~$\{x_i, \boldsymbol{\eta}_i\}$ follows a stationary and ergodic process;
(2)~$\mathbf{g}_i$ is a martingale difference sequence with
$\E(\mathbf{g}_i \mathbf{g}_i') = \mathbf{S}$, which is a positive definite matrix;
(3)~$\E(x_i^2) = \sigma_x^2 > 0$;
(4)~$\beta \neq 0$.

Let $\hat{\delta} \equiv s_{xz}^{-1} s_{xy}$ where
\[
  s_{xz} \equiv \frac{1}{n} \sum_{i=1}^{n} x_i z_i, \quad
  s_{xy} \equiv \frac{1}{n} \sum_{i=1}^{n} x_i y_i.
\]

\begin{enumerate}[label=(\alph*)]
\item Prove that $\sigma_{xz} \equiv \E(x_i z_i) \neq 0$ (so that the rank condition for
  identification is satisfied).

\item Prove that $\hat{\delta} \pto \delta$.
\end{enumerate}

We now want to consider a case in which the rank condition is just barely
satisfied, that is, $\sigma_{xz}$ is very close to zero.  One way to do this is to replace
Assumption~(4) with

\noindent (4$'$) $\beta = 1/\sqrt{n}$.

For the remainder of the problem, assume that (1)--(3) and (4$'$) hold.  (This
assumption and the following questions are based on Staiger and Stock (1997).)
Note that $\hat{\delta} - \delta = s_{xz}^{-1} \bar{g}_1$, where $\bar{g}_1$ is the sample mean
of $g_{1i}$ $(\equiv x_i \varepsilon_i)$.

\begin{enumerate}[label=(\alph*),resume]
\item Show that $s_{xz} \pto 0$.

\item Show that $\sqrt{n}\, s_{xz} \dto \sigma_x^2 + a$, where $a$ is distributed
  $N(0, s_{22})$ and $s_{22}$ is the $(2,2)$ element of $\mathbf{S}$.

\item Show that $\hat{\delta} - \delta \dto (\sigma_x^2 + a)^{-1} b$, where $(a, b)$ are
  jointly normally distributed with mean zero and covariance matrix $\mathbf{S}$.

\item Is $\hat{\delta}$ consistent?
\end{enumerate}

\bigskip
\noindent\textsc{Empirical Exercises}
\bigskip

Read Griliches (1976) before answering.  We will estimate the type of the wage
equation estimated by Griliches using an extract from the NLS-Y used by Blackburn
and Neumark (1992).\footnote{See Blackburn and Neumark (1992, Section~III) for a more
detailed description of the sample.}  The NLS-Y is \textbf{panel data}, with the same set of
young men surveyed at several points in time (we will not exploit the panel feature of the
data set in this exercise, though).  The extract contains information about those
individuals at two points in time: first, the earliest year in which wages and other
variables are available, and second, in 1980.  In a data file GRILIC.ASC, data are
provided on $RNS$, $RNS80$, $MRT$, $MRT80$, $SMSA$, $SMSA80$, $MED$, $IQ$, $KWW$,
$YEAR$, $S$, $S80$, $EXPR$, $EXPR80$, $TENURE$, $TENURE80$, $LW$, and $LW80$ (in
this order, with the columns corresponding to the variables, as usual).  The variable
definitions are:
\begin{align*}
  RNS &= \text{dummy for residency in the southern states} \\
  MRT &= \text{dummy for marital status (1 if married)} \\
  SMSA &= \text{dummy for residency in metropolitan areas} \\
  MED &= \text{mother's education in years} \\
  KWW &= \text{score on the ``Knowledge of the World of Work'' test} \\
  IQ &= \text{IQ score} \\
  AGE &= \text{age of the individual} \\
  S &= \text{completed years of schooling} \\
  EXPR &= \text{experience in years} \\
  TENURE &= \text{tenure in years} \\
  LW &= \text{log wage.}
\end{align*}
$YEAR$ is the year of the first point in time.  Variables without ``80'' are for the first
point, and those with ``80'' are for 1980.

The Blackburn--Neumark sample has 815 observations after deleting black individuals.
Also deleted are cases with missing information on mother's education,
reducing the sample size to 758.

\begin{enumerate}[label=(\alph*)]
\item Calculate means and standard deviations of all the provided variables (including
  those for 1980) and prepare a table similar to Griliches's Table~1.  Also,
  calculate the correlation between $IQ$ and $S$.

\item Consider the wage equation (dropping the individual subscript $i$)
  \[
    LW = \beta S + \gamma\, IQ + \boldsymbol{\delta}'\mathbf{h} + \varepsilon,
  \]
  where $LW$ is log wages, $S$ is schooling, and
  \[
    \mathbf{h} \equiv (EXPR, TENURE, RNS, SMSA, \text{year dummies})'.
  \]
  (There is no need to include a constant because it is a linear combination of
  the year dummies.)  Prepare a table similar to Table~3.2.  In your 2SLS estimation
  of the wage equation, the set of instruments should consist of predetermined
  regressors ($S$ and $\mathbf{h}$) and excluded predetermined variables ($MED$,
  $KWW$, $MRT$, $AGE$).  Is the relative magnitude of the three different estimates
  of the schooling coefficient consistent with the prediction based on the omitted
  variable and errors-in-variables biases?

\item For your 2SLS estimation of the wage equation in part~(b), calculate Sargan's
  statistic (it should be 87.655).  What should the degrees of freedom be?  Calculate
  the $p$-value.

\item Obtain the 2SLS estimate by actually running two regressions.  Verify that the
  standard errors given by the second stage regression are \emph{different} from those
  you obtained in~(b).

\item Griliches mentions that schooling, too, may be endogenous.  What is his argument?
  Estimate by 2SLS the wage equation, treating both $IQ$ and $S$ as endogenous.
  What happens to the schooling coefficient?  How would you explain the
  difference between your 2SLS estimate of the schooling coefficient here and
  your 2SLS estimate in~(b)?  Calculate Sargan's statistic (it should be 13.268)
  and its $p$-value.

\item Estimate the wage equation by GMM, treating schooling as predetermined as in
  the 2SLS estimation in part~(b).  Are 2SLS standard errors smaller than GMM
  standard errors?  Test whether schooling is predetermined by the $C$ statistic.
  To calculate $C$, you need to do two GMMs with and without schooling as an
  instrument.  Use the same $\hat{\mathbf{S}}$ throughout.  To calculate $\hat{\mathbf{S}}$, use the 2SLS
  residual from part~(b).  The $C$ statistic should be 58.168.

\item Go back to the wage equation you estimated in~(e) by 2SLS endogenous schooling.
  The large Sargan's statistic (and the large $J$ statistic in part~(f)) is a concern.
  So consider dropping both $MED$ and $KWW$ from the list of instruments.
  Is the order condition for identification still satisfied?  What happens to the
  2SLS estimates?  (The schooling coefficient will be $-529.3$\%!)  What do you
  think is the problem?  \textbf{Hint:} There are two endogenous regressors, $S$ and $IQ$.
  $MRT$ and $AGE$ are the only excluded predetermined variables.  If $MRT$ and $AGE$
  do not have explanatory power in the first-stage regression of $S$ and $IQ$
  on the instruments, then the fitted values of $S$ and $IQ$ will be very close to linear
  combinations of the included predetermined variables.  This will lead to a near-multicollinearity
  problem in the second-stage regression.  Make sure you examine
  the first-stage regressions to check the explanatory power of $MRT$ and $AGE$.
\end{enumerate}

\end{problemset}

\begin{answers}

\noindent\textsc{Analytical Exercises}
\bigskip

\answer{4}
The asymptotic variance of the OLS estimator is (by setting $\mathbf{x} = \mathbf{z}$ in (3.5.1))
\[
  \boldsymbol{\Sigma}_{\mathbf{xz}}^{-1}\, \mathbf{A}\,
  \boldsymbol{\Sigma}_{\mathbf{xz}}^{-1}
  = (\boldsymbol{\Sigma}_{\mathbf{zz}}\, \mathbf{A}^{-1}\,
  \boldsymbol{\Sigma}_{\mathbf{zz}})^{-1},
\]
where $\mathbf{A} \equiv \E(\varepsilon_i^2\, \mathbf{z}_i \mathbf{z}_i')$.  The asymptotic variance
of the two-step efficient GMM estimator is, by (3.5.13),
\[
  (\boldsymbol{\Sigma}_{\mathbf{xz}}' \mathbf{S}^{-1}
  \boldsymbol{\Sigma}_{\mathbf{xz}})^{-1},
\]
where $\mathbf{S} \equiv \E(\varepsilon_i^2\, \mathbf{x}_i \mathbf{x}_i')$.  If $\mathbf{W}$ is as
defined in the hint, then
\[
  \mathbf{W}\mathbf{S}\mathbf{W} = \mathbf{W} \quad \text{and} \quad
  \boldsymbol{\Sigma}_{\mathbf{xz}}' \mathbf{W} \boldsymbol{\Sigma}_{\mathbf{xz}}
  = \boldsymbol{\Sigma}_{\mathbf{zz}} \mathbf{A}^{-1} \boldsymbol{\Sigma}_{\mathbf{zz}}.
\]
So (3.5.1) reduces to the asymptotic variance of the OLS estimator.  By (3.5.11),
it is no smaller than
$(\boldsymbol{\Sigma}_{\mathbf{xz}}' \mathbf{S}^{-1} \boldsymbol{\Sigma}_{\mathbf{xz}})^{-1}$,
which is the asymptotic variance of the two-step efficient GMM estimator.

\bigskip
\noindent\textsc{Empirical Exercises}
\bigskip

\answer{(b)}
See Table~3.3, lines~1--5.

\answer{(e)}
He gives three reasons for the endogeneity of schooling.  First, the measurement
error $\eta$ may be related to success in schooling.  The second reason is the
argument we made in the text: the choice of $S$ is influenced by unobservable
characteristics.  Third, schooling may be measured with error.  The general
formula for the asymptotic bias for the 2SLS estimator is given by
\[
  \plim\, \hat{\boldsymbol{\delta}}_{\text{2SLS}} - \boldsymbol{\delta}
  = (\boldsymbol{\Sigma}_{\mathbf{xz}}' \boldsymbol{\Sigma}_{\mathbf{xx}}^{-1}
  \boldsymbol{\Sigma}_{\mathbf{xz}})^{-1}\,
  \boldsymbol{\Sigma}_{\mathbf{xz}}' \boldsymbol{\Sigma}_{\mathbf{xx}}^{-1}\,
  \E(\mathbf{x}_i \cdot \varepsilon_i).
\]
If $S$ is endogenous, the instrument set $\mathbf{x}_i$ for the 2SLS estimation in~(b)
erroneously includes a variable (which is $S$) that is not predetermined.  So one of the
elements of $\E(\mathbf{x}_i \cdot \varepsilon_i)$ is not zero.  Then it is clear from the formula that
the asymptotic bias can exist at least for some coefficients.

\answer{(g)}
Without $MED$ and $KWW$ as instruments, the 2SLS estimator becomes imprecise.
The order condition is satisfied with equality.  The culprit is the low
explanatory power of $MRT$ (look at its $t$-value in its first-stage regressions for
$S$ and $IQ$).
\end{answers}
